\chapter{Fuentes y alcance}
Las prácticas y explicaciones son originales. Los ejemplos se contrastan con el código de tinygrad fijado al commit c3aec477b99d9bb87c54d91897cf60acd3f17441. La documentación web puede cambiar; ante una diferencia de API, prevalece el fuente de esa revisión para reproducir esta edición.

\begin{thebibliography}{9}
\bibitem{tiny-revision} Tiny Corp y colaboradores. \emph{tinygrad}, revisión del 30-09-2026. \url{https://github.com/tinygrad/tinygrad/tree/c3aec477b99d9bb87c54d91897cf60acd3f17441}.
\bibitem{tiny-tensor} tinygrad. \emph{Tensor}: constructor, realización y backward. \url{https://github.com/tinygrad/tinygrad/blob/c3aec477b99d9bb87c54d91897cf60acd3f17441/tinygrad/tensor.py}.
\bibitem{tiny-optim} tinygrad. Optimizadores y contexto de entrenamiento. \url{https://github.com/tinygrad/tinygrad/blob/c3aec477b99d9bb87c54d91897cf60acd3f17441/tinygrad/nn/optim.py}.
\bibitem{tiny-jit} tinygrad. Captura y reutilización de TinyJit. \url{https://github.com/tinygrad/tinygrad/blob/c3aec477b99d9bb87c54d91897cf60acd3f17441/tinygrad/engine/jit.py}.
\bibitem{tiny-device} tinygrad. Selección de dispositivo, renderer e interfaz en helpers. \url{https://github.com/tinygrad/tinygrad/blob/c3aec477b99d9bb87c54d91897cf60acd3f17441/tinygrad/helpers.py}.
\bibitem{tiny-docs} tinygrad. Documentación oficial y guía del runtime. \url{https://docs.tinygrad.org/}. \url{https://docs.tinygrad.org/developer/runtime/}.
\end{thebibliography}
